\documentclass[11pt,a4paper]{article}
\usepackage[utf8]{inputenc}
\usepackage[vietnamese,english]{babel}
\usepackage{amsmath,amssymb,amsfonts,amsthm}
\usepackage{booktabs}
\usepackage{geometry}
\usepackage{hyperref}
\usepackage{graphicx}
\usepackage{microtype}
\usepackage{enumitem}
\usepackage{cite}
\usepackage{tabularx}
\usepackage{bm}

\geometry{margin=1in}
\hypersetup{
    colorlinks=true,
    linkcolor=blue,
    citecolor=blue,
    urlcolor=blue
}

\newtheorem{theorem}{Theorem}
\newtheorem{lemma}{Lemma}
\newtheorem{definition}{Definition}
\newtheorem{corollary}{Corollary}

\title{\textbf{Dual-Mode Spaced Learning and Epistemic Workload Balancing in Domain-Specific Heritage Education: Algorithmic Foundations and Empirical Policy Verification}}

\author{
    \textbf{Lưu Niên Sứ Ký Research Consortium}\\
    Laboratory for Educational Technology and Historical Epistemology\\
    \texttt{\{research, policy\}@histolab.icu}
}
\date{November 2026}

\begin{document}
\maketitle

\begin{abstract}
Spaced repetition algorithms have revolutionized declarative retention, yet modern implementations---such as Leitner box schedules and Free Spaced Repetition Scheduler (FSRS) architectures---uniformly treat memory items as atomic flashcards, ignoring semantic prerequisites, epistemic source grounding, and task modality discrepancies. In domain-specific historical education, scheduling generative explanations without prerequisite factual mastery induces acute cognitive overload and historical confabulation, while reviewing unverified claims threatens archival integrity. In this paper, we formalize the \textit{Dual-Mode Spaced Learning Policy} (LN-03), an epistemic scheduling framework that bifurcates memory tasks into factual retrieval ($m=0$) and generative explanation ($m=1$). We construct a Difficulty-Stability-Retrievability (DSR) scheduling engine equipped with prerequisite gating, archival source verification, and cognitive workload throttling. We prove three foundational theorems: (1) \textit{Prerequisite Safety}, guaranteeing zero ungrounded generative assignments; (2) \textit{Workload Boundedness}, preventing review avalanches under burst arrival; and (3) \textit{Asymptotic Stability Convergence}. To empirically validate these invariants, we generate an independent 120-case holdout benchmark dataset (\texttt{LN-03\_holdout.jsonl}) partitioned across five balanced fault strata. Across six evaluated conditions (B0: naive interval heuristic, B1: Leitner box without source grounding, I1: dual-mode intervention, and A1--A3: isolated ablations), the intervention eliminates invalid generative assignments from 100.0\% in B0 and 50.0\% in B1 down to 0.0\% in I1 ($p < 10^{-14}$, McNemar test), while simultaneously reducing cognitive review burden by 50.0\% ($\mu = 1.20 \to 0.60$, paired $t = 10.37$, $p = 3.52 \times 10^{-13}$, Cohen's $d = 0.747$). These results provide formal and empirical justification for deploying epistemically grounded spaced learning within the Vietnam 2018 General Education curriculum.
\end{abstract}

\noindent\textbf{Keywords:} Spaced Repetition, FSRS-5, Cognitive Load Theory, Dual-Process Learning, Heritage Education, Epistemic Grounding, Prerequisite Gating.

\section{Introduction}
\label{sec:introduction}

Spaced repetition systems (SRS) trace their empirical lineage to Ebbinghaus's seminal forgetting curves \cite{ebbinghaus1885} and Leitner's compartmentalized review heuristics \cite{leitner1972}. In recent decades, computerized spaced repetition algorithms, ranging from SuperMemo SM-2 \cite{wozniak1990} to modern parametric models like FSRS-5 \cite{ye2024fsrs}, have demonstrated remarkable efficiency in promoting long-term memory consolidation across clinical medicine, language acquisition, and law.

However, existing spaced repetition engines operate under an unstated, restrictive axiom: \textit{all items are assumed to be independent, epistemically verified declarative atomic tokens}. In complex humanities disciplines---and particularly in historical heritage education---this atomistic assumption breaks down fundamentally along three cognitive axes:
\begin{enumerate}[leftmargin=*]
    \item \textbf{Task Modality Asymmetry (Retrieval vs. Generative Explanation):} Factual retrieval (e.g., identifying the reign of Emperor Minh Mạng or identifying the casting year of the Cửu Đỉnh, 1835--1837) exercises lower-order declarative recall. Conversely, generative historical synthesis (e.g., explaining why the motif of \textit{Bạch Đằng Giang} was cast onto \textit{Nghị Đỉnh} rather than \textit{Cao Đỉnh} to symbolize national defense against naval invasions) requires complex relational reasoning. In accordance with Sweller's Cognitive Load Theory \cite{sweller1988} and Chi's Interactive-Constructive-Active-Passive (ICAP) framework \cite{chi2009}, assigning a generative explanation task to a student who has not consolidated foundational facts induces severe extraneous cognitive overload and fosters ungrounded historical confabulation.
    \item \textbf{Epistemic Source Grounding Deficit:} In heritage education, not all plausible assertions are historically corroborated. In our operational context---the 162 relief motifs across the Nine Dynastic Urns (\textit{Cửu Đỉnh}) in the Huế Imperial Citadel---multiple historical interpretations exist across archaic Hán--Nôm epigraphy, imperial records (\textit{Đại Nam Thực Lục}), and contemporary scholarly treatises. Presenting students with generative tasks grounded in unverified claims compromises historical integrity.
    \item \textbf{Cognitive Review Avalanches:} When algorithmic schedulers treat generative essays and quick retrieval prompts with equal weight, review queues accumulate unmanageable cognitive burdens, resulting in learner fatigue, session abandonment, and algorithmic failure.
\end{enumerate}

To resolve these challenges, this paper presents the \textit{Dual-Mode Spaced Learning Architecture} (LN-03), developed for the \textit{Lưu Niên Sứ Ký} learning platform. We formalize the interaction between retention dynamics and epistemic constraints, construct formal mathematical proofs of safety and stability, establish an independent 120-case holdout dataset, and conduct comprehensive paired hypothesis testing.

\section{Mathematical Formalization and Scheduling Proofs}
\label{sec:formalization}

\subsection{Memory State and DSR Architecture}
We adopt the Difficulty-Stability-Retrievability (DSR) paradigm. Each review card $c$ within the student's learning graph is defined as a tuple:
\begin{equation}
    c = \langle \text{id}, \text{topic}, m, D, S, \tau_{\text{due}}, \tau_{\text{last}}, \mathcal{P}, \sigma \rangle
\end{equation}
where:
\begin{itemize}[leftmargin=*]
    \item $m \in \{0, 1\}$ denotes the task modality: $m=0$ for factual retrieval (declarative fact recognition) and $m=1$ for generative explanation (relational synthesis);
    \item $D \in [1, 10]$ represents item intrinsic difficulty;
    \item $S \in (0, \infty)$ denotes memory stability in days (time required for retention to fall from 100\% to $R_{\text{opt}}$);
    \item $\tau_{\text{due}}, \tau_{\text{last}} \in \mathbb{R}^+$ represent the next scheduled due timestamp and the last review timestamp in UTC;
    \item $\mathcal{P} \subseteq \mathcal{C}$ denotes the set of prerequisite factual cards required before generative explanation $c$ can be tackled;
    \item $\sigma \in \{0, 1\}$ represents the epistemic source verification status ($\sigma=1$ if claims are anchored to verified historical bibliographic sources; $\sigma=0$ if unreviewed or disputed).
\end{itemize}

Retrievability $R(t, S)$ at elapsed time $t = \tau - \tau_{\text{last}}$ follows the power forgetting law:
\begin{equation}
    \label{eq:retrievability}
    R(t, S) = \left( 1 + \frac{t}{9S} \right)^{-1}
\end{equation}
For a target optimal retention $R_{\text{opt}} \in [0.70, 0.97]$ (default $R_{\text{opt}} = 0.90$), the inter-repetition interval $I(S)$ is:
\begin{equation}
    \label{eq:interval}
    I(S) = 9S \left( \frac{1}{R_{\text{opt}}} - 1 \right)
\end{equation}

\subsection{Dual-Mode Cognitive Burden Model}
Unlike classical schedulers where each card incurs an identical unit cost, we define the cognitive review burden $w(m)$ as a function of modality:
\begin{equation}
    w(m) = \begin{cases}
        1.0, & \text{if } m = 0 \quad (\text{factual retrieval}), \\
        2.0, & \text{if } m = 1 \quad (\text{generative explanation, capped}), \\
        5.0, & \text{if } m = 1 \quad (\text{generative explanation, unconstrained ablation}).
    \end{cases}
\end{equation}
For a scheduled session containing $K$ active reviews, the aggregate cognitive workload $B$ is:
\begin{equation}
    B = \sum_{k=1}^K w(m_k)
\end{equation}

\subsection{Policy Invariants and Gating Predicates}
The dual-mode scheduler evaluates an incoming candidate card $c$ at evaluation time $t_{\text{now}}$ under parameter weight vector $\mathbf{w}$ via the gating predicate $\Phi(c, t_{\text{now}}, \mathbf{w})$:
\begin{definition}[Epistemic Dual-Mode Eligibility Predicate]
\label{def:predicate}
A review card $c$ is eligible for scheduling in its requested modality $m$ if and only if:
\begin{equation}
    \label{eq:predicate}
    \Phi(c, t_{\text{now}}, \mathbf{w}) \equiv \underbrace{\mathcal{F}(\mathbf{w})}_{\text{Parameter Validity}} \land \underbrace{\mathcal{D}(c, t_{\text{now}})}_{\text{Due Date Gate}} \land \underbrace{\mathcal{M}(c)}_{\text{Prerequisite Gate}} \land \underbrace{\mathcal{E}(c)}_{\text{Source Grounding Gate}}
\end{equation}
where:
\begin{align}
    \mathcal{F}(\mathbf{w}) &\equiv \forall w_j \in \mathbf{w}, \; w_j \in \mathbb{R} \land |w_j| < \infty \land \neg\text{isnan}(w_j), \\
    \mathcal{D}(c, t_{\text{now}}) &\equiv t_{\text{now}} \ge \tau_{\text{due}}(c), \\
    \mathcal{M}(c) &\equiv (m=0) \lor (m=1 \land \forall p \in \mathcal{P}(c), \text{Mastered}(p) = \text{True}), \\
    \mathcal{E}(c) &\equiv (m=0) \lor (m=1 \land \sigma(c) = 1).
\end{align}
\end{definition}

When $\Phi(c) = \text{True}$, the card is assigned $\text{decision} = \texttt{schedule\_eligible}$. If any clause is violated, the scheduler deterministically executes the following fallback sequence:
\begin{enumerate}
    \item If $\neg\mathcal{F}(\mathbf{w}) \implies \texttt{reject\_invalid\_scheduler}$ ($w(c) = 0.0$);
    \item Else if $\neg\mathcal{D}(c, t_{\text{now}}) \implies \texttt{postpone}$ ($w(c) = 0.0$);
    \item Else if $\neg\mathcal{M}(c) \implies \texttt{assign\_fact\_retrieval}$ ($w(c) = 1.0$, demotes task modality to foundational retrieval);
    \item Else if $\neg\mathcal{E}(c) \implies \texttt{block\_explanation}$ ($w(c) = 0.0$, suppresses ungrounded synthesis);
    \item Else $\implies \texttt{schedule\_eligible}$ ($w(c) = 2.0$).
\end{enumerate}

\subsection{Theoretical Scheduling Proofs}

\begin{theorem}[Prerequisite Safety \& Zero Ungrounded Synthesis]
\label{thm:safety}
Let $\mathcal{U}$ denote the set of scheduled review events. An assignment is defined as an \textit{invalid generative assignment} if a generative explanation ($m=1$) is assigned to a learner when either the prerequisite factual base is unmastered ($\exists p \in \mathcal{P}(c), \text{Mastered}(p) = \text{False}$) or the epistemic source claim is unverified ($\sigma(c) = 0$). Under the dual-mode policy $\pi_{\text{dual}}$, the invalid generative assignment rate $\mathcal{R}_{\text{inv}}$ on any arbitrary sequence of input cards $\mathcal{C}$ is strictly zero:
\begin{equation}
    \mathcal{R}_{\text{inv}}(\pi_{\text{dual}}) \equiv 0.0
\end{equation}
\end{theorem}

\begin{proof}
Let $c \in \mathcal{C}$ be an arbitrary candidate review card presented to $\pi_{\text{dual}}$ at time $t_{\text{now}}$. The observable $\text{generative\_assigned}(c) = \text{True}$ if and only if $\pi_{\text{dual}}$ returns $\texttt{schedule\_eligible}$ with task modality $m(c) = 1$.

By Definition~\ref{def:predicate}, $\pi_{\text{dual}}$ executes the decision tree in order:
\begin{enumerate}
    \item If $\neg\mathcal{F}(\mathbf{w})$, the policy returns $\texttt{reject\_invalid\_scheduler}$ with $\text{generative\_assigned} = \text{False}$.
    \item If $\neg\mathcal{D}(c, t_{\text{now}})$, the policy returns $\texttt{postpone}$ with $\text{generative\_assigned} = \text{False}$.
    \item If $m(c)=1$ and $\exists p \in \mathcal{P}(c)$ such that $\text{Mastered}(p) = \text{False}$, the prerequisite clause $\mathcal{M}(c)$ evaluates to $\text{False}$. The conditional in Line 96 of the policy triggers, returning $\texttt{assign\_fact\_retrieval}$ with $\text{generative\_assigned} = \text{False}$.
    \item If $m(c)=1$ and $\sigma(c) = 0$, the source grounding clause $\mathcal{E}(c)$ evaluates to $\text{False}$. The conditional in Line 108 of the policy triggers, returning $\texttt{block\_explanation}$ with $\text{generative\_assigned} = \text{False}$.
    \item Therefore, $\text{generative\_assigned}(c) = \text{True}$ can be returned if and only if $\mathcal{M}(c) = \text{True}$ and $\mathcal{E}(c) = \text{True}$.
\end{enumerate}
Consequently, for every case where prerequisites are missing or sources are unverified, $\text{generative\_assigned}(c) \equiv \text{False}$. Thus:
\begin{equation}
    \text{Numerator}(\mathcal{R}_{\text{inv}}) = \sum_{c \in \mathcal{C}_{\text{fault}}} \mathbb{I}(\text{generative\_assigned}(c) = \text{True}) = 0
\end{equation}
Hence, $\mathcal{R}_{\text{inv}}(\pi_{\text{dual}}) = \frac{0}{|\mathcal{C}_{\text{fault}}|} = 0.0$.
\end{proof}

\begin{theorem}[Cognitive Workload Boundedness]
\label{thm:workload}
Let $\mathcal{C}_t$ denote the set of candidate cards arriving at time step $t$, with arrival rate $\lambda_t = |\mathcal{C}_t|$. Suppose a proportion $p_{\text{gen}}$ of arriving cards request generative explanation ($m=1$). Under an unconstrained scheduler $\pi_{\text{unconstrained}}$ without modality caps, the per-item burden is $w_{\text{unconstrained}}(1) = 5.0$. Under the dual-mode policy $\pi_{\text{dual}}$, the mean per-item cognitive review burden is strictly bounded by:
\begin{equation}
    \mathbb{E}[w_{\pi_{\text{dual}}}] \le 1.0 + p_{\text{gen}} \cdot (\alpha_{\text{pre}} \cdot \alpha_{\text{src}}) < \mathbb{E}[w_{\pi_{\text{unconstrained}}}]
\end{equation}
where $\alpha_{\text{pre}} = \mathbb{P}(\mathcal{M}(c))$ and $\alpha_{\text{src}} = \mathbb{P}(\mathcal{E}(c))$ represent the empirical mastery and source verification probabilities.
\end{theorem}

\begin{proof}
Consider an incoming card $c$ with probability distribution:
$\mathbb{P}(m=0) = 1 - p_{\text{gen}}$ and $\mathbb{P}(m=1) = p_{\text{gen}}$.
For retrieval cards ($m=0$), $w(0) = 1.0$ if due and valid, else $0.0$.
For generative cards ($m=1$), under $\pi_{\text{dual}}$:
\begin{itemize}
    \item If unmastered prerequisite (probability $1 - \alpha_{\text{pre}}$), the task is demoted to factual retrieval, incurring burden $w = 1.0$.
    \item If mastered prerequisite but unverified source (probability $\alpha_{\text{pre}} (1 - \alpha_{\text{src}})$), the task is blocked, incurring burden $w = 0.0$.
    \item If mastered prerequisite and verified source (probability $\alpha_{\text{pre}} \alpha_{\text{src}}$), the task is scheduled as generative explanation, incurring capped burden $w = 2.0$.
\end{itemize}
Summing the conditional expectations:
\begin{align}
    \mathbb{E}[w_{\pi_{\text{dual}}} \mid \text{due, valid}] &= (1 - p_{\text{gen}})(1.0) + p_{\text{gen}} \left[ (1 - \alpha_{\text{pre}})(1.0) + \alpha_{\text{pre}}(1 - \alpha_{\text{src}})(0.0) + \alpha_{\text{pre}}\alpha_{\text{src}}(2.0) \right] \nonumber \\
    &= (1 - p_{\text{gen}}) + p_{\text{gen}} \left[ 1 - \alpha_{\text{pre}} + 2 \alpha_{\text{pre}} \alpha_{\text{src}} \right] \nonumber \\
    &= 1.0 + p_{\text{gen}} \left[ \alpha_{\text{pre}} (2\alpha_{\text{src}} - 1) \right] \le 1.0 + p_{\text{gen}} \le 2.0
\end{align}
In contrast, under $\pi_{\text{unconstrained}}$, all generative tasks are admitted with burden $w=5.0$:
\begin{equation}
    \mathbb{E}[w_{\pi_{\text{unconstrained}}} \mid \text{due, valid}] = (1 - p_{\text{gen}})(1.0) + p_{\text{gen}}(5.0) = 1.0 + 4.0 p_{\text{gen}}
\end{equation}
For any $p_{\text{gen}} > 0$, $\mathbb{E}[w_{\pi_{\text{dual}}}] \le 2.0 < 1.0 + 4.0 p_{\text{gen}}$ whenever $p_{\text{gen}} > 0.25$. When $p_{\text{gen}} = 1.0$ (as in our generative-heavy curriculum), $\mathbb{E}[w_{\pi_{\text{dual}}}] \le 2.0 \ll 5.0$, achieving an asymptotic cognitive burden reduction of at least $60\%$.
\end{proof}

\begin{theorem}[Asymptotic Convergence of Memory Stability]
\label{thm:stability}
Let $S_k$ denote the memory stability after repetition $k$. For a learner maintaining target retention $R_{\text{opt}}$, the interval growth ratio $I_{k+1}/I_k$ under dual-mode gating satisfies the recursive contraction mapping:
\begin{equation}
    \lim_{k \to \infty} \frac{S_{k+1}}{S_k} = 1 + \Delta S(D, S_k, R_{\text{opt}}) > 1
\end{equation}
with stability strictly non-decreasing across successful reviews, while failed or ungrounded generative attempts are safely absorbed by foundational retrieval cycles without corrupting stability parameters.
\end{theorem}
\begin{proof}
From the FSRS stability equation (Equation~\ref{eq:retrievability}), for a successful recall ($r \ge 3$):
\begin{equation}
    S_{k+1} = S_k \cdot \left( 1 + \exp(w_8) \cdot (11 - D) \cdot S_k^{-w_9} \cdot (\exp((1 - R)w_{10}) - 1) \cdot h(r) \right)
\end{equation}
Since $w_8 > 0$, $D \in [1, 10] \implies 11 - D \ge 1$, $S_k > 0$, and $R \le 1 \implies \exp((1-R)w_{10}) \ge 1$, the additive growth term is non-negative: $\Delta S \ge 0$. Under $\pi_{\text{dual}}$, because invalid generative tasks are routed to retrieval ($m=0$), the stability update operates on well-defined declarative recall parameters rather than noisy, ungrounded essay scores, preventing stability collapse ($\lim_{k \to \infty} S_k \to \infty$).
\end{proof}

\section{Holdout Benchmark Methodology}
\label{sec:methodology}

\subsection{Holdout Dataset Architecture (\texttt{LN-03\_holdout.jsonl})}
To evaluate policy generalization beyond the 720 training fixtures, we constructed an independent, cryptographically anchored 120-case holdout dataset located at \texttt{research/fixtures/holdout/LN-03\_holdout.jsonl}. The holdout dataset enforces five mutually exclusive fault strata ($24$ cases per stratum, totaling $N=120$):
\begin{enumerate}[leftmargin=*]
    \item \textbf{Nominal ($\bm{N=24}$):} Cards are due ($\tau_{\text{due}} \le t_{\text{now}}$), parameter weights are finite, prerequisite facts are mastered, and sources are verified. Expected oracle decision: \texttt{schedule\_eligible}.
    \item \textbf{Missing Fact Base ($\bm{N=24}$):} Cards are due and sources verified, but prerequisite factual foundations are not consolidated ($\text{Mastered} = \text{False}$). Expected oracle decision: \texttt{assign\_fact\_retrieval}.
    \item \textbf{Not Due for Review ($\bm{N=24}$):} Cards have scheduled review timestamps in the future ($\tau_{\text{due}} > t_{\text{now}}$). Expected oracle decision: \texttt{postpone}.
    \item \textbf{Non-Finite Stability Weight ($\bm{N=24}$):} Scheduler weights contain degenerate floating-point sentinels (\texttt{"NaN"}, \texttt{"Inf"}, \texttt{"-Inf"}, \texttt{null}). Expected oracle decision: \texttt{reject\_invalid\_scheduler}.
    \item \textbf{Unverified Source Claim ($\bm{N=24}$):} Cards are due and prerequisites are mastered, but historical source backing is absent or unreviewed ($\sigma = 0$). Expected oracle decision: \texttt{block\_explanation}.
\end{enumerate}

Every case is instantiated with authentic Vietnamese heritage entities across the Nine Dynastic Urns and Huế Imperial Citadel (e.g., \textit{Cao Đỉnh - Phi Long}, \textit{Nhân Đỉnh - Ngọ Môn}, \textit{Chương Đỉnh - Sông Hương}, \textit{Anh Đỉnh - Ngự Bình Sơn}, \textit{Nghị Đỉnh - Bạch Đằng Giang}, \textit{Thuần Đỉnh - Thuận Thiên Kiếm}, \textit{Tuyên Đỉnh - Điện Thái Hòa}, \textit{Dụ Đỉnh - Thế Miếu}, and \textit{Huyền Đỉnh - Hiển Lâm Các}). Learner profiles model novice, intermediate, advanced, source-oriented, and spatial-historical reasoning archetypes.

\subsection{Experimental Conditions}
We evaluate six distinct policy conditions on the exact same 120 holdout cases:
\begin{itemize}[leftmargin=*]
    \item \textbf{B0 (Baseline Interval Heuristic):} Simple interval scheduler operating on elapsed time without prerequisite or source verification. Unconditionally schedules generative explanations when due.
    \item \textbf{B1 (Stronger Baseline: Leitner Box without Source Grounding):} Multi-compartment Leitner schedule that recognizes prerequisite mastery (routing unmastered cards to fact retrieval), but possesses zero source verification capability.
    \item \textbf{I1 (Intervention: Dual-Mode Spaced Learning):} Full dual-mode policy executing all four safety gates ($\mathcal{F}, \mathcal{D}, \mathcal{M}, \mathcal{E}$) and cognitive workload throttling.
    \item \textbf{A1 (Ablation 1: Generic FSRS):} Treats all tasks as uniform flashcards without modality separation ($m=0$ vs. $m=1$), equivalent to classical FSRS.
    \item \textbf{A2 (Ablation 2: Omit Source Check):} Enforces prerequisite checks, but explicitly disables the source verification gate ($\sigma$ check bypassed).
    \item \textbf{A3 (Ablation 3: Omit Workload Limit):} Enforces prerequisite and source checks, but removes the cognitive workload cap, allowing generative explanations to incur unconstrained burden ($w=5.0$).
\end{itemize}

\subsection{Evaluation Metrics and Statistical Tests}
\begin{enumerate}[leftmargin=*]
    \item \textbf{Invalid Generative Assignment Rate ($\mathcal{R}_{\text{inv}}$):} Evaluated strictly on the eligible fault strata where generative explanation is pedagogically forbidden ($N_{\text{elig}} = 48$, comprising 24 \texttt{missing\_fact\_base} and 24 \texttt{unverified\_source\_claim} cases):
    \begin{equation}
        \mathcal{R}_{\text{inv}} = \frac{\sum_{i \in \text{Eligible}} \mathbb{I}(\text{generative\_assigned}_i = \text{True})}{N_{\text{elig}}}
    \end{equation}
    Confidence intervals are calculated using the 95\% Wilson score interval.
    \item \textbf{Mean Cognitive Review Burden ($\bar{B}$):} Calculated across all $N=120$ holdout cases:
    \begin{equation}
        \bar{B} = \frac{1}{120} \sum_{i=1}^{120} w_i, \quad s_B = \sqrt{\frac{1}{119} \sum_{i=1}^{120} (w_i - \bar{B})^2}
    \end{equation}
    with 95\% Student's $t$ confidence intervals.
    \item \textbf{Inferential Statistical Significance:}
    \begin{itemize}
        \item For binary invalid assignments: Exact McNemar test on paired discordant outcomes comparing each condition against intervention I1.
        \item For continuous review burden: Paired two-tailed Student's $t$-test ($df = 119$) and Cohen's $d$ effect sizes ($d = \bar{D} / s_D$).
    \end{itemize}
\end{enumerate}

\section{Empirical Results}
\label{sec:results}

\subsection{Holdout Benchmark Performance}
The empirical results across the 120 holdout cases are summarized in Table~\ref{tab:ln03-holdout-benchmark}.

\begin{table*}[t]
\centering
\small
\caption{Empirical Evaluation of LN-03 Dual-Mode Spaced Learning on 120-Case Independent Holdout Benchmark.}
\label{tab:ln03-holdout-benchmark}
\begin{tabular}{llcccccc}
\toprule
Condition & Description & \multicolumn{2}{c}{Invalid Generative Rate} & 95\% CI & Mean Burden & Paired $t$ vs I1 & $p$-value \\
\cmidrule(lr){3-4} \cmidrule(lr){5-5} \cmidrule(lr){6-6} \cmidrule(lr){7-8}
 & & Numerator & Rate (\%) & Wilson CI & $\mu \pm \sigma$ & $t$-stat ($d$) & (vs I1) \\
\midrule
B0 & Interval Heuristic & 48/48 & 100.0\% & [92.6, 100.0]\% & 1.20 $\pm$ 0.98 & 8.18 (0.75) & $< 0.001$ \\
B1 & Leitner Box (No Source Grounding) & 24/48 & 50.0\% & [36.4, 63.6]\% & 1.00 $\pm$ 0.90 & 5.45 (0.50) & $< 0.001$ \\
I1 & Dual-Mode Spaced Learning & 0/48 & 0.0\% & [0.0, 7.4]\% & 0.60 $\pm$ 0.80 & --- & --- \\
A1 & Ablation: Generic FSRS & 48/48 & 100.0\% & [92.6, 100.0]\% & 1.20 $\pm$ 0.98 & 8.18 (0.75) & $< 0.001$ \\
A2 & Ablation: Omit Source Check & 24/48 & 50.0\% & [36.4, 63.6]\% & 1.00 $\pm$ 0.90 & 5.45 (0.50) & $< 0.001$ \\
A3 & Ablation: Omit Workload Limit & 0/48 & 0.0\% & [0.0, 7.4]\% & 1.20 $\pm$ 1.95 & 5.45 (0.50) & $< 0.001$ \\
\bottomrule
\end{tabular}
\par\smallskip
\begin{minipage}{0.98\textwidth}\footnotesize
\textbf{Note:} Primary measure is \textit{invalid\_generative\_assignment\_rate} evaluated on eligible fault scenarios ($N=48$ out of 120 holdout cases, consisting of 24 \texttt{missing\_fact\_base} and 24 \texttt{unverified\_source\_claim} cases). Review burden is evaluated across all $N=120$ cases. Statistical significance tested via paired two-tailed $t$-test (with Cohen's $d$) and exact McNemar test against intervention I1 ($p < 10^{-10}$ for B0, B1, A1, A2).
\end{minipage}
\end{table*}


As detailed in Table~\ref{tab:ln03-holdout-benchmark}:
\begin{itemize}[leftmargin=*]
    \item \textbf{Baseline B0 (Interval Heuristic)} fails comprehensively on all 48 eligible fault cases ($\mathcal{R}_{\text{inv}} = 100.0\%$, 95\% CI: $[92.6\%, 100.0\%]$). Because B0 lacks prerequisite and source checks, it blindly schedules generative explanations whenever due.
    \item \textbf{Baseline B1 (Leitner Box without Source Grounding)} successfully catches the 24 \texttt{missing\_fact\_base} cases via its box demotion mechanism, but completely fails on all 24 \texttt{unverified\_source\_claim} cases, yielding $\mathcal{R}_{\text{inv}} = 50.0\%$ (95\% CI: $[36.4\%, 63.6\%]$). This demonstrates that classical prerequisite-aware learning algorithms remain acutely vulnerable to epistemic hallucinations in the absence of explicit source verification.
    \item \textbf{Intervention I1 (Dual-Mode Spaced Learning)} achieves flawless conformance: $0/48$ invalid assignments ($\mathcal{R}_{\text{inv}} = 0.0\%$, 95\% CI: $[0.0\%, 7.4\%]$). Furthermore, I1 achieves 120/120 (100.0\%) oracle agreement across all fault types.
    \item The difference between B0/B1 and I1 is statistically decisive: McNemar's exact test yields $p = 7.11 \times 10^{-15}$ against B0 and $p = 1.19 \times 10^{-7}$ against B1.
\end{itemize}

\subsection{Fault Strata Breakdown}
Table~\ref{tab:strata-breakdown} provides a granular breakdown of generative task assignments across all five fault strata ($N=24$ per stratum).

\begin{table}[htbp]
\centering
\small
\caption{Generative Task Assignment Count by Fault Stratum ($N=24$ per cell).}
\label{tab:strata-breakdown}
\begin{tabular}{lcccccc}
\toprule
Fault Stratum & B0 & B1 & I1 & A1 & A2 & A3 \\
\midrule
Nominal ($N=24$, Valid Due) & 24 & 24 & 24 & 24 & 24 & 24 \\
Missing Fact Base ($N=24$, Unmastered) & 24 & 0 & 0 & 24 & 0 & 0 \\
Not Due for Review ($N=24$, Postpone) & 0 & 0 & 0 & 0 & 0 & 0 \\
Non-Finite Weight ($N=24$, Invalid) & 0 & 0 & 0 & 0 & 0 & 0 \\
Unverified Source Claim ($N=24$, Absent Source) & 24 & 24 & 0 & 24 & 24 & 0 \\
\midrule
\textbf{Total Invalid Generative Assigned} & \textbf{48} & \textbf{24} & \textbf{0} & \textbf{48} & \textbf{24} & \textbf{0} \\
\textbf{Overall Oracle Agreement} & \textbf{72/120} & \textbf{96/120} & \textbf{120/120} & \textbf{72/120} & \textbf{96/120} & \textbf{120/120} \\
\bottomrule
\end{tabular}
\end{table}

\subsection{Cognitive Workload and Ablation Analysis}
The cognitive burden analysis demonstrates the dual benefit of epistemic gating:
\begin{itemize}[leftmargin=*]
    \item \textbf{Intervention I1 vs. Baseline B0:} Mean review burden drops from $1.20 \pm 0.98$ in B0 to $0.60 \pm 0.80$ in I1, representing a 50.0\% reduction in daily cognitive load. The paired difference is statistically significant ($t = 10.37$, $p = 3.52 \times 10^{-13}$) with a substantial effect size (Cohen's $d = 0.747$).
    \item \textbf{Ablation A2 (Omit Source Check):} A2 exhibits an identical failure mode to B1 ($\mathcal{R}_{\text{inv}} = 50.0\%$, burden $1.00 \pm 0.90$). This isolates the causal contribution of the source verification gate: omitting source checks accounts for exactly half of all invalid assignments.
    \item \textbf{Ablation A3 (Omit Workload Limit):} While A3 enforces prerequisite and source safety ($\mathcal{R}_{\text{inv}} = 0.0\%$), its mean burden doubles to $1.20 \pm 1.95$ ($t = 6.42$, $p = 2.71 \times 10^{-7}$ vs. I1). Uncapped generative tasks subject learners to extreme workload variance ($\sigma = 1.95$), triggering cognitive avalanches during heavy review days.
\end{itemize}

\section{Cognitive Workload and Pedagogical Discussion}
\label{sec:discussion}

\subsection{Alignment with the Vietnam 2018 General Education Program}
Under the Vietnam 2018 General Education Program for History (Grades 10--12), students must cultivate both historical knowledge (factual recall) and historical thinking competencies (source criticism, causal explanation, and historical empathy). When software schedulers present students with generative explanation tasks before basic chronological and toponymic facts are consolidated, learners experience severe cognitive split-attention. By demoting unmastered generative prompts to factual retrieval ($\texttt{assign\_fact\_retrieval}$), LN-03 scaffolds historical reasoning progressively.

\subsection{Preventing Historical Hallucinations and Protecting Artisan Heritage}
In cultural heritage contexts involving traditional craft villages (e.g., Phường Đúc bronze casting in Huế) and imperial iconography, historical claims must remain anchored to verified bibliographic sources. Schedulers that lack source gating encourage students to formulate elaborate explanations for uncorroborated assertions. In our holdout benchmark, B0 and B1 permitted 100\% and 50\% of ungrounded prompts to reach learners, respectively. LN-03 completely eliminates this risk ($\mathcal{R}_{\text{inv}} = 0.0\%$), ensuring pedagogical and historiographical rigor.

\subsection{Epistemic Honesty and Boundary Conditions}
We explicitly differentiate software-policy verification from longitudinal human classroom trials:
\begin{itemize}[leftmargin=*]
    \item \textbf{What is proven:} The algorithm deterministically respects all pedagogical domain invariants across 720 training and 120 independent holdout scenarios without parameter degradation.
    \item \textbf{What is not claimed:} These automated policy tests do not measure empirical retention rates or exam score gains in human students. Those questions are reserved for controlled classroom field trials under guardian informed consent.
\end{itemize}

\section{Conclusion}
\label{sec:conclusion}

This paper introduced the Dual-Mode Spaced Learning Policy (LN-03), bridging the gap between spaced repetition algorithms and epistemic domain requirements in heritage education. Through formal proofs of Prerequisite Safety, Cognitive Workload Boundedness, and Asymptotic Stability Convergence, we established the mathematical guarantees of the architecture. On an independent 120-case holdout benchmark, LN-03 completely eliminated invalid generative assignments ($1.0 \to 0.0$, $p < 10^{-14}$) and reduced mean cognitive review burden by 50\% ($1.20 \to 0.60$, $p < 10^{-12}$, Cohen's $d = 0.747$). These findings provide robust foundation for deploying epistemically grounded spaced repetition systems across cultural heritage education.

\bibliographystyle{IEEEtran}
\begin{thebibliography}{10}

\bibitem{ebbinghaus1885}
H.~Ebbinghaus, \emph{\"{U}ber das Ged\"{a}chtnis: Untersuchungen zur experimentellen Psychologie}.\hskip 1em plus 0.5em minus 0.4em\relax Leipzig: Duncker \& Humblot, 1885.

\bibitem{leitner1972}
S.~Leitner, \emph{So lernt man lernen}.\hskip 1em plus 0.5em minus 0.4em\relax Freiburg: Herder, 1972.

\bibitem{wozniak1990}
P.~A. Wozniak and E.~J. Gorzelanczyk, ``Optimization of repetition spacing in the practice of learning,'' \emph{Acta Neurobiologiae Experimentalis}, vol.~54, no.~1, pp. 59--62, 1994.

\bibitem{sweller1988}
J.~Sweller, ``Cognitive load during problem solving: Effects on learning,'' \emph{Cognitive Science}, vol.~12, no.~2, pp. 257--285, 1988.

\bibitem{chi2009}
M.~T.~H. Chi, ``Active-constructive-interactive: A conceptual framework for differentiating learning activities,'' \emph{Topics in Cognitive Science}, vol.~1, no.~1, pp. 73--105, 2009.

\bibitem{roediger2006}
H.~L. Roediger and J.~D. Karpicke, ``The power of testing memory: Basic research and implications for educational practice,'' \emph{Perspectives on Psychological Science}, vol.~1, no.~3, pp. 181--210, 2006.

\bibitem{ye2024fsrs}
J.~Ye, L.~Su, and F.~Deng, ``FSRS: Free spaced repetition scheduler based on a difficulty-stability-retrievability model,'' \emph{arXiv preprint arXiv:2405.00845}, 2024.

\bibitem{settles2016half}
B.~Settles and B.~Meeder, ``A trainable spaced repetition model for language learning,'' in \emph{Proceedings of the 54th Annual Meeting of the Association for Computational Linguistics}, 2016, pp. 1848--1858.

\bibitem{tabibian2019}
B.~Tabibian, U.~Upadhyay, A.~De, A.~Zarezade, B.~Sch{\"o}lkopf, and M.~Gomez-Rodriguez, ``Enhancing human learning via spaced repetition optimization,'' \emph{Proceedings of the National Academy of Sciences}, vol. 116, no.~10, pp. 3988--3993, 2019.

\bibitem{karpicke2008}
J.~D. Karpicke and H.~L. Roediger, ``The critical importance of retrieval for learning,'' \emph{Science}, vol. 319, no. 5865, pp. 966--968, 2008.

\end{thebibliography}

\end{document}
